\chapter{Introduction} \label{cap:intro}

\section{Motivation}

\textit{Computer Organization and Design: RISC-V Edition} by Patterson and Hennessy \cite{patterson2020}
is one of the most widely adopted references in undergraduate computer architecture courses worldwide.
At the University of São Paulo's Escola Politécnica, the course \textit{PCS3225 --- Digital Systems 2}
uses this textbook as its primary reference. Students follow its progression step by step, implementing
a simplified RISC-V processor in Verilog from book reference.

The simplified single-cycle processor introduced in Chapter 4.4 of the book, \textit{A Simple
Implementation Scheme}, supports only eight instructions: \texttt{lw}, \texttt{sw}, \texttt{beq},
\texttt{add}, \texttt{addi}, \texttt{sub}, \texttt{and}, and \texttt{or}. This restriction is
deliberate and pedagogically motivated, it isolates the essential datapath and control concepts before
more complex features are introduced. A subsequent homework assignment extends the processor with
\texttt{lui} and \texttt{jalr}, enabling function calls. Both implementations are partial subsets of
the RV32I base ISA \cite{riscv-spec}.

This restriction creates a practical barrier. The standard RISC-V GCC toolchain
(\texttt{riscv32-unknown-elf-gcc}) targets the full RV32I base ISA \cite{riscv-spec}, which includes
shift instructions, byte and halfword memory operations, multiple branch variants, and PC-relative
addressing. Any C program compiled with this toolchain will emit instructions that the student-built
processor cannot execute.

As a result, students are limited to running hand-written assembly or small programs provided by the
professor. They cannot compile and run arbitrary C code on the processor they designed and built. This
prevents them from connecting the hardware they implement to the software abstractions they use daily,
and forecloses the pedagogical opportunity of observing how a compiler translates and optimizes C code
into the primitive instructions their processor supports.

\section{Objective}

This work develops eight progressive GCC compiler targets for the simplified RISC-V processors
described in the Hennessy and Patterson textbook, enabling students of the PCS3225 course at USP to
compile and run C programs on the processors they build in Verilog.

The specific objectives are:

\begin{enumerate}
  \item Define a minimal GCC target (\texttt{rvsc0}) matching the eight-instruction single-cycle
        processor described in Chapter 4.4 of the textbook, synthesizing all operations not natively
        supported by that processor.
  \item Define a target (\texttt{rvsc1}) matching the course homework processor, \texttt{rvsc0}
        extended with \texttt{lui} and \texttt{jalr}, as the minimum instruction set capable of
        supporting the full C calling convention with synthesis.
  \item Define six additional progressive targets (\texttt{rvsc2} through \texttt{rvsc7}),
        incrementally extending the supported instruction set from bare RV32I through RV64IMAFD,
        for students who wish to continue developing their processor beyond the course scope.
  \item Synthesize every instruction not natively supported by a given target as an equivalent
        sequence of instructions that the target does support.
\end{enumerate}

\section{Rationale}

To the best of the author's knowledge, no prior work addresses C compilation for intentionally
restricted pedagogical RISC-V subsets. GCC's machine description framework \cite{gcc-internals} makes
it possible to define a backend that synthesizes missing instructions transparently from the primitives
the hardware does support. Applying this mechanism to pedagogical ISAs that deliberately omit standard
instructions appears to be novel.

The ability to run a self-written C program on a processor the student designed and built closes a
pedagogical loop that rarely closes in undergraduate education. Most courses treat hardware and software
as adjacent subjects that never directly intersect. This work creates a complete vertical slice from C
source to register-level execution on student hardware. Students can compile a function, inspect the
output, and observe concretely how the compiler synthesizes a shift operation from repeated additions,
or a signed comparison from subtraction and bit manipulation, making the cost of each ISA restriction
tangible rather than abstract.

Furthermore, students interact with GCC, the dominant open-source compiler for embedded and systems
software, rather than a pedagogical toy. The flags, ABI conventions, ELF output, and linker scripts
they encounter are identical to those used in professional and research settings, giving the exercise
relevance beyond the course itself.

\section{Document Organization}

The remainder of this monograph is organized as follows.
Chapter~\ref{cap:related-work} surveys related work, in instruction synthesis for embedded compiler
backends and in tools for instruction-set-level computing education.
Chapter~\ref{cap:background} presents the conceptual background: the RISC-V instruction set, the
single-cycle datapath of the Hennessy-Patterson educational processor, the C calling convention the
targets must honour, and the structure of a GCC backend.
Chapter~\ref{cap:method} describes the development method, including the derivation cycle applied to
each synthesized instruction and the validation approach.
Chapter~\ref{cap:requirements} specifies the requirements for each of the eight targets, defining
the allowed instruction sets and the synthesis obligations imposed on the compiler.
Chapter~\ref{cap:development} details the implementation: a correctness argument and an instruction
sequence for every synthesis, the machine-description and per-target configuration that make them
transparent to the programmer, and the limitations that remain.
Chapter~\ref{cap:results} reports the three test layers, quantifies the static and dynamic cost of
synthesis on the Embench-IoT suite, and discusses what the measurements do and do not establish.
Chapter~\ref{cap:conclusion} presents the conclusions, contributions, and directions for future work.
The appendices collect the compiler output that the synthesis sections refer to, including a worked
compilation example across rvsc0, rvsc1 and rvsc2.
